%!TEX root = forallxyyc.tex

\ifHTMLtarget\else
% Título de rosto

\pagestyle{empty}

\vspace*{80pt}

\begin{raggedleft}
\fontsize{30pt}{24pt}\sffamily
\selectfont
  \textbf{forall 
  {\fontsize{37pt}{24pt}\selectfont\rmfamily\textit{x}}: 
  Calgary}

\medskip\fontsize{18pt}{20pt}\selectfont

\textbf{Uma introdução à\\ lógica formal}

\vfill
\fontsize{12pt}{16pt}\selectfont \textit{Por } \textbf{P.~D. Magnus}\\
\textbf{Tim Button}\\
\textbf{Robert Trueman}\\
\textbf{Richard Zach}\\
\textit{com contribuições de}\\
\textbf{J.~Robert Loftis}\\
\textbf{Aaron Thomas-Bolduc}\\

\medskip
\textbf{Tradução para o português brasileiro}\\
Carlos André Duarte Costa\\
26 de setembro de 2026

\vfill
\textbf{\forallxversion}\par
\end{raggedleft}

\newpage

\noindent\small%
\fi
\ifHTMLtarget
\par\bigskip\noindent
\textbf{Tradução para o português brasileiro}\par
Carlos André Duarte Costa\par
26 de setembro de 2026\par
\bigskip
\fi
Este livro é baseado em
\href{https://www.homepages.ucl.ac.uk/~uctytbu/OERs.html}{\forallx:
\textit{Cambridge}}, de
\href{https://www.homepages.ucl.ac.uk/~uctytbu/}{Tim Button} (University College London),
utilizado sob uma licença \href{https://creativecommons.org/licenses/by/4.0/}{CC BY
4.0}, que por sua vez se baseia
em \href{https://www.fecundity.com/logic/}{\forallx}, de
\href{https://www.fecundity.com/job/}{P.D.\ Magnus}
(University at Albany, State University of New York),
utilizado sob uma licença \href{https://creativecommons.org/licenses/by/4.0/}{CC BY
4.0},
e foi remixado, revisado e ampliado por Aaron
Thomas-Bolduc e \href{https://richardzach.org/}{Richard Zach}
(University of Calgary).
Ele inclui material adicional de \forallx{} por P.~D. Magnus e
\href{https://www.homepages.ucl.ac.uk/~uctytbu/OERs.html}{\textit{Metatheory}}, de Tim Button,
utilizado sob uma licença \href{https://creativecommons.org/licenses/by/4.0/}{CC BY
4.0},
de \href{https://github.com/rob-helpy-chalk/openintroduction}{\forallx: \textit{Lorain
County Remix}},
de \href{https://sites.google.com/site/cathalwoods/}{Cathal Woods} e
J. Robert Loftis, e de \href{http://www.rtrueman.com/uploads/7/0/3/2/70324387/modal_logic_primer.pdf}{\textit{A Modal Logic Primer}}, de \href{http://www.rtrueman.com/}{Robert Trueman}, reproduzido com permissão específica.
Esse material não está abrangido pela licença CC BY 4.0 mencionada abaixo.

\par\bigskip
\noindent\textbf{Revisão editorial da edição brasileira.}
Equipe de revisão da edição brasileira.

\ifHTMLtarget
\chapter*{Créditos do projeto REALMat}
\else
\newpage
\thispagestyle{empty}
\begin{center}
{\Large\sf\textbf{Créditos do projeto REALMat}}
\end{center}
\fi
\noindent
Projeto REALMat — Recursos Educacionais Abertos na Licenciatura em Matemática.

\medskip
\noindent
O REALMat organiza e disponibiliza esta edição brasileira.

\smallskip
\noindent
Portal: \href{https://projetorealmat.github.io/forallx-yyc/}{projetorealmat.github.io/forallx-yyc}.

\noindent
Repositório desta edição: \href{https://github.com/projetorealmat/forallx-yyc}{github.com/projetorealmat/forallx-yyc}.

\noindent
Fonte original: \href{https://github.com/rzach/forallx-yyc}{github.com/rzach/forallx-yyc}.

\ifHTMLtarget\else
  \bigskip

  \noindent\footnotesize
\fi
Os materiais da edição disponibilizados sob a licença \href{https://creativecommons.org/licenses/by/4.0/}{Creative Commons Atribuição 4.0} são regidos por essa licença. O material de \textit{A Modal Logic Primer} permanece reproduzido com permissão específica e não é oferecido sob CC BY 4.0.
Você pode copiar e redistribuir o material em qualquer meio ou formato, bem como remixá-lo, transformá-lo e criar obras derivadas para qualquer finalidade, inclusive comercialmente, nos termos seguintes:
\begin{itemize}
\item Você deve atribuir o devido crédito, fornecer um link para a licença e indicar se foram feitas alterações. Isso pode ser feito de qualquer maneira razoável, mas não de modo que sugira que o licenciador apoia você ou o uso que você faz da obra.
\item Você não pode aplicar termos jurídicos nem medidas tecnológicas que restrinjam legalmente outras pessoas de fazer qualquer coisa que a licença permita.
\end{itemize}

\ifHTMLtarget\else
\vfil\normalsize\noindent
\fi
O código-fonte \LaTeX{} deste livro está disponível no GitHub em
\href{https://github.com/rzach/forallx-yyc/}{github.com/rzach/forallx-yyc/} e nos formatos PDF e
HTML em
\href{https://forallx.openlogicproject.org}{forallx.openlogicproject.org}.
Esta versão é a revisão \gitAbbrevHash{}
(\gitCommitterDate).

\bigskip
\noindent
A preparação deste livro-texto foi possível graças a uma bolsa do \href{https://taylorinstitute.ucalgary.ca/}{Taylor Institute for Teaching and Learning}.

\bigskip
\noindent
\href{https://taylorinstitute.ucalgary.ca/}{%
%  \iflatexml
    \includegraphics[alt={Logotipo do Taylor Institute for Teaching and Learning}]
      {assets/ti-color}
  % \else
  %   \includegraphics[width=8cm]{assets/ti-color}
  % \fi
  }

\ifHTMLtarget
O formato HTML é produzido com
\href{https://vlmantova.github.io/bookml/}{BookML}, de
\href{https://eps.leeds.ac.uk/maths/staff/4058/dr-vincenzo-l-mantova}{Vincenzo Mantova},
e \href{https://dlmf.nist.gov/LaTeXML/}{\LaTeX{}ML}, de
\href{https://www.nist.gov/people/bruce-r-miller}{Bruce Miller} e
\href{https://prodg.org/}{Deyan Ginev}.
\else
\bigskip
\noindent Projeto da capa: Mark Lyall.
\fi
